\documentclass[../RFC-HDFG-2026-001.tex]{subfiles}

\begin{document}

\begin{abstract}

    
This RFC proposes a string-based configuration API for HDF5 I/O filters.
The current \texttt{H5Pset\_filter} interface requires users to construct raw
\texttt{cd\_values} arrays of \texttt{unsigned int}, forcing error-prone type-punning for
floating-point parameters and producing opaque configurations that are
hostile to CLI tools and high-level language bindings.

The proposed design
introduces \texttt{H5Pappend\_filter}, which accepts a filter ID and
either a raw \texttt{cd\_values} array or a TOML inline table parameter
string via a tagged union \texttt{H5Z\_params\_t}, resolving string parameters
into standard \texttt{cd\_values} at configuration time. The parameter string
format is a subset of the TOML v1.0.0 inline table syntax~\cite{toml10},
parsed via a vendored subset of the tomlc17 library; this provides
typed values (integer, float, boolean, string), nested configuration
for future VOL and VFD use, and a stable externally-maintained grammar.
Filters opt into string configuration via a new \texttt{H5Z\_class3\_t} plugin class with
optional \texttt{set\_config}/\texttt{get\_config} callbacks; the
existing \texttt{name} field's contract is strengthened from
free-form debug comment to a required canonical name that serves as
both the stable string identifier and the display name.  The library
writes this name into the filter-pipeline message at filter-add time
so it is available to tools even when the plugin is not installed.
Existing call sites compile and function without modification; files written by a v3-aware library are readable by older libraries through v2 plugins that tolerate extra trailing \texttt{cd\_values} slots (see \S\ref{sec:class3-struct}).
\end{abstract}

\vfill
{\footnotesize\textbf{Acknowledgements:} This work is supported by the U.S.\
Department of Energy, Office of Science, Advanced Scientific Computing Research,
under Contract DE-AC02-06CH11357 and by the U.S.\ Department of Energy, Office
of Science, Office of Advanced Scientific Computing Research, Scientific
Discovery through Advanced Computing (SciDAC) program.}

\end{document}
